\begin{recipe}{Boerenkoolstamp met rookworst}
\kicker[Hoofdgerecht,Vlees,Nederlands,Doordeweeks]{Hoofdgerecht · vlees · Nederlands · doordeweeks}
\index[register]{Boerenkoolstamp met rookworst}
\index[register]{Boerenkool}
\index[register]{Rookworst}
\index[register]{Stamppot}

\meta{35 minuten}

\lettrine{B}{oerenkoolstamp} met een Gelderse rookworst is de klassieke
winterstamp. In mijn jeugd aten we er een plak roggebrood bij, en soms wat
zuurkool, al lust ik die laatste tegenwoordig als enige nog.

\blockrule{Ingrediënten}
\begin{ingredients}
  \ing{600 g}{boerenkool (zak), schoongemaakt}\index[register]{Boerenkool}
  \ing{1}{Gelderse rookworst (375 g)}\index[register]{Rookworst}
  \ing{1 kg}{kruimige aardappelen, geschild en in stukken}\index[register]{Aardappelen}
  \ing{1 el}{mosterd}
  \ingb{scheutje melk}
  \ingb{klontje boter}
  \ingb{versgemalen peper en zout}
\end{ingredients}

\blockrule{Bereiding}
\tip{Laat het water niet koken terwijl de rookworst erin ligt. Te heet
water laat de worst barsten of taai worden.}
\begin{steps}
  \item Kook de aardappelen in gezouten water gaar, ongeveer 20 minuten.
  \item Warm de rookworst ondertussen op in niet-kokend water, volgens de
        aanwijzingen op de verpakking.
  \item Giet de aardappelen af en laat goed uitstomen. Stamp ze grof, met
        nog wat stukjes erin, samen met een klein scheutje melk, een
        klontje boter en de mosterd. Voeg liever te weinig melk toe dan te
        veel: je kunt er altijd meer bij doen, terug halen is lastiger.
  \item Roer de boerenkool in delen door de hete puree, tot hij geslonken
        is. De puree wordt vanzelf gladder doordat de boerenkool erdoorheen
        mengt, dus stamp hem daarvoor niet te fijn.
  \item Breng op smaak met peper en zout. Snijd de rookworst in plakjes
        en verdeel ze over de borden. Zet een pot mosterd op tafel voor
        wie meer wil.
\end{steps}
\end{recipe}
